\chapter{Results}
\label{cap:resultados}

\section{Data gate}

All 56 primary files were present. Twenty-eight partitions were generated. BTCUSDT
produced 887,803 events; ETHUSDT, 939,101. The independent window produced 1,288,352
events. The availability and volume gate was therefore satisfied.

\section{Reproduction at the fixed point}

\input{en/e2_fixed_point.tex}

Table~\ref{tab:e2-ponto-fixo} reproduces the qualitative pattern of \citet{lee2023}.
The MLE is practically unbiased, and its standard deviations are about \(\sqrt{2}\)
times those in the paper, as expected for paths of half the length. The networks show
the characteristic shrinkage bias: they pull \(\alpha_{\mathrm{self}}\) and
\(\alpha_{\mathrm{cross}}\) toward the center of the prior, as in the paper, in
exchange for lower variance. The two versions of the MLE produced exactly the same
estimates, so the initialization at the true value used by Lee does not explain the
MLE's advantage at this well-identified point (\(g\approx0.73\)). The macro MSE was
0.0103 for the MLE and 0.0111 for the LSTM ensemble.

\section{Final comparison: LSTM, Transformer and MLE}

\input{en/final_results.tex}

\input{en/posthoc_results.tex}

\section{Computational cost}
\label{sec:custo}

The final run recorded 105~ms of wall-clock time per path for the MLE with four
processes. This value was measured while neural training was occupying the GPU. A
controlled benchmark on the same machine, with the same code (which exactly
reproduces the 15,000 recorded estimates), produced the values in
Table~\ref{tab:custo}.

\begin{table}[htbp]
\centering
\caption{Cost per path of 1,000 events in a controlled benchmark.}
\label{tab:custo}
\small
\begin{tabular}{lr}
\toprule
Estimator and execution mode & Time per path\\
\midrule
MLE, 5 starts, one core & 9.6 ms\\
MLE, 1 start, one core & 2.2 ms\\
MLE, 5 starts, 4 processes (wall-clock) & 2.5 ms\\
LSTM, batch size one, CPU (1 thread) / MPS & 7.0 ms / 29.6 ms\\
Transformer, batch size one, CPU (1 thread / 10 threads) / MPS & 9.9 / 4.8 / 5.0 ms\\
LSTM, batch of 256 on MPS & 0.07 ms\\
Transformer, batch of 32 on MPS & 3.7 ms\\
\midrule
\citet{lee2023}: network / MLE & 12 ms / 1,763 ms\\
\bottomrule
\end{tabular}
\end{table}

With an analytic gradient and a linear recursion, the MLE of this four-parameter
model costs a few milliseconds, the same order as neural inference at batch size one.
The ratio of about 150 times reported by Lee is not reproduced. The cost advantage of
the networks appears in batched inference, especially for the LSTM, and would grow
for models with more expensive likelihoods.

Lee does not describe the implementation of the MLE. To find where its 1,763~ms come from,
the same paths at Lee's fixed point were fitted with implementations that differ in one factor
at a time, at 1,000 and 2,000 events, the two lengths the paper mentions
(Appendix~\ref{ap:lee}). All variants reach the same estimates (largest difference 0.002 in any
parameter); only the cost changes (Table~\ref{tab:custo-variantes}). The likelihood is the
dominant factor: computed as a double sum over pairs of events, it costs about 90 times as much
as the recursion at 2,000 events, while a numerical gradient and a derivative-free optimizer
(Nelder--Mead) multiply the cost by about five each. With the recursion, even an interpreted
loop stays below 0.1~s, and the time reported by Lee falls between the two vectorized
double-sum variants at 2,000 events. The reported ratio is therefore most likely due to a
likelihood computed without the recursion of Section~\ref{sec:mle-theory}.

\begin{table}[htbp]
\centering
\caption{Cost of one MLE fit by implementation, at the fixed point of \citet{lee2023}: median of
10 paths, one start at the true parameters, one core. The pure-Python double sum was timed for
one evaluation and multiplied by the number of evaluations of the vectorized Nelder--Mead fit.}
\label{tab:custo-variantes}
\small
\begin{tabular}{llrr}
\toprule
Likelihood & Gradient and optimizer & 1,000 events & 2,000 events\\
\midrule
Recursion, vectorized & analytic, L-BFGS-B (this work) & 1.1 ms & 1.7 ms\\
Recursion, vectorized & numerical, L-BFGS-B & 5.4 ms & 8.3 ms\\
Recursion, Python loop & numerical, L-BFGS-B & 47 ms & 89 ms\\
Double sum, vectorized & numerical, L-BFGS-B & 176 ms & 711 ms\\
Double sum, vectorized & Nelder--Mead & 818 ms & 3,358 ms\\
Double sum, Python loop & Nelder--Mead (extrapolated) & 26 s & 105 s\\
\midrule
\multicolumn{3}{l}{\citet{lee2023}, implementation not described} & 1,763 ms\\
\bottomrule
\end{tabular}
\end{table}

\section{Empirical calibration}

The MLE converged on all 1,063 calibration windows. The details are in
Section~\ref{sec:calibracao}: cross-excitation is practically zero, \(\beta\) exceeds
the upper limit of the prior in most windows, and no window falls inside all the
ranges used in training. The following sections report the empirical protocol of
Section~\ref{sec:empirical-protocol}.

\input{empirical_results.tex}

\input{addendum_results.tex}
